@@M0@@
\uline{where we put}, \uline{for} @@M1@@:
@@M2@@
\uline{and for} @@M3@@ \uline{and} @@M4@@:
@@M5@@
\uline{with the notation of} 3.1 \uline{and} 3.3, \uline{and} @@M6@@ \uline{denoting the} @@M7@@-\uline{st Chern class of} @@M8@@ (V 6.7). \uline{Moreover}, @@M9@@ \uline{has a left inverse} @@M10@@, \uline{defined by} @@M11@@.

\uline{If in addition} @@M12@@ \uline{is noetherian and} @@M13@@ \uline{regular}, \uline{we have the exact sequence}:
@@M14@@

\emph{Proof.} Observe that, when @@M15@@ is regular, one could give a purely algebraic proof in terms of the theory of @@M16@@-rings, using the description of the situation in 4.7. We shall instead follow the proof of 3.7.

We know that the ``inverse-image'' homomorphisms @@M17@@, being @@M18@@-homomorphisms, are compatible with the @@M19@@-ring filtrations, and consequently define homomorphisms on the associated gradeds; hence the homomorphisms:
@@M20@@
and the analogous homomorphisms obtained after tensoring with @@M21@@. Furthermore, 4.6 shows that after tensoring with @@M22@@, the homomorphisms @@M23@@ pass to the associated gradeds, and therefore define:
@@M24@@
